\documentclass[11pt,letterpaper]{article}
\usepackage[utf8]{inputenc}
\usepackage[T1]{fontenc}
\usepackage{geometry}
\geometry{a4paper, margin=2cm}
\usepackage{hyperref}
\usepackage{setspace}
\usepackage{siunitx}
\sisetup{separate-uncertainty=true, range-phrase=--}
\DeclareSIUnit\kcalmol{kcal\,mol^{-1}}
\DeclareSIUnit\angstrom{\text{\AA}}
\singlespacing

\begin{document}
\emergencystretch=2em

\begin{flushleft}
\textbf{Myke-Vital Temgoua Sao} \\
Department of Physics, Faculty of Science \\
University of Yaound\'{e} I \\
P.O. Box 812, Yaound\'{e}, Cameroon \\
\href{mailto:myke-vital.sao@facsciences-uy1.cm}{myke-vital.sao@facsciences-uy1.cm} \\[1.5em]
August 28, 2026
\end{flushleft}

\begin{flushleft}
Prof.\ Kenneth M.\ Merz, Jr., Editor-in-Chief \\
\textit{Journal of Chemical Information and Modeling} \\
American Chemical Society
\end{flushleft}

\vspace{0.5em}

\noindent Dear Prof.\ Merz,

We submit our manuscript \textbf{``Resistance-Aware Docking Stratifies Antimalarial Leads by Mutant Retention but Short MD Shows Non-Equivalent Estimands''} for consideration as an Article in \textit{JCIM}.

The central question is whether resistance-aware polypharmacology docking can prioritize antimalarial leads that retain activity across clinically relevant mutant panels, and whether short molecular dynamics simulations corroborate or refute that prioritization. Seventeen Set-C candidates were ranked by a docking-derived resistance-retention score (RRS) across six PfDHFR and PfCRT resistance states, with chemical-space (ACSI) and network (PNS) scores kept as separate ranking dimensions. The primary analysis on 12 two-target candidates yielded a discriminant five-class stratification (A*: 1, A: 1, B: 4, C: 5, D: 1) that separates predicted wild-type potency from predicted mutant retention. Four parent-study complexes that are not on the Set-C list were simulated for \qty{10}{\nano\second} each as a short pose-retention check; only PfCRT--214 gave an interpretable MM-GBSA estimate ($\qty{-18.25 \pm 0.40}{\kcalmol}$). A 16-system single-replicate MD pilot on two Set-C candidates did not recover the docking retention pattern; trajectory metrics and docking-RRS disagreed in direction for seven of eight mutant comparisons, demonstrating that docking-RRS and MD-RRS are non-equivalent estimands rather than interchangeable measures of the same quantity. An external docking panel (39 ligands, 312 records) averaged near RRS 100. The ranking is informative inside this computational panel only as a hypothesis-generating prioritization, not as a validated triage procedure; short MD stress tests probe a different estimand and cannot serve as direct validation of docking-derived retention classes.

Code, source data, and analysis scripts are available on GitHub under an open-source licence; a Zenodo snapshot will be deposited at publication. Statistics are reported with effect sizes, permutation $p$-values, bootstrap intervals, and Bonferroni adjustment.

\noindent\textbf{Suggested reviewers.}
\begin{enumerate}
  \item Dr.\ Fidele Ntie-Kang, University of Buea --- African natural-product computational discovery
  \item Dr.\ Kelly Chibale, University of Cape Town --- antimalarial drug discovery
  \item Dr.\ John D.\ Chodera, Memorial Sloan Kettering --- free-energy methods and MD best practices
\end{enumerate}

\noindent This manuscript is original, not under consideration elsewhere, and all authors have approved the submission.

\vspace{1em}

\noindent Sincerely,

\vspace{1em}

\noindent \textbf{Myke-Vital Temgoua Sao} \\
(on behalf of all co-authors)

\end{document}
